\documentclass{article}
\usepackage[
    paperwidth=6in,
    paperheight=8.5in,
    hmargin=0in,
    vmargin=0in,
    nohead,
    nofoot,
    includeheadfoot
]{geometry}
\setlength{\topskip}{0mm}
\setlength{\parindent}{0mm}
\usepackage{tikz}
\usepackage{scalefnt}
\usetikzlibrary{positioning, calc}
\usepackage{datetime}
\usepackage{xstring}
\usepackage{pgfmath}
\def\hauteurLigneJourSemaine{0.225}



\tikzset{
    circle node/.style = {circle,inner sep=1pt,draw, fill=white, minimum size=3pt},
    X node/.style = {fill=white, inner sep=1pt},
    day node/.style = {font=\scalefont{0.7}, yshift=-3.5pt},
    task node/.style = { font=\scalefont{1}},
    event node/.style = {circle, fill=red},
    dot node/.style = {circle, draw, inner sep=5pt, }
}

\newcommand{\monthTaskDot}[1]{
    \fill (#1) circle [radius=1pt];
}
 

\begin{document}
\setlength{\parindent}{0pt} %Retirer les marges.
\thispagestyle{empty} % Disable page numbers
\thispagestyle{empty}
    


    % Create dotted paper background
    \begin{tikzpicture}[remember picture, overlay]
        %Séparer la page en deux 

        
        %Dessiner les bullets points : 
        \def\dotRadius{0.6pt}  % Size of the dots
        \newcounter{lineNum}   % Define a counter for line numbering
        \setcounter{lineNum}{0} %
        \foreach \x in {1.5, 1.6875, ..., 6} {
             %\setcounter{lineNum}{0}
             %\node[anchor=east] at (\x in, 0 in) {\thelineNum};
             %\stepcounter{lineNum}
            
            \foreach \y in {0.230, 0.46, ..., 8.5} {
                %\node[anchor=east] at (0 in, -\y in) {\thelineNum};
                \fill[gray!70] (\x in, -\y  in) circle (\dotRadius);
                %\stepcounter{lineNum}
            }
        }

        \foreach \y in {\hauteurLigneJourSemaine, 1.604, ..., 8.5} {
            \draw[gray!90] (1.5in, -\y  in) to (6in, -\y in);
        }

        \node[anchor=north west, font=\bfseries] at (1.95in, 0in) {Lundi};
        \node[anchor=north west, font=\bfseries] at (3.45in, 0in) {Mardi};
        \node[anchor=north west, font=\bfseries] at (4.9in, 0in) {Mercredi};
        %Espace verticale vide 
        \draw [gray!90] (1.5in, 0 in) to (1.5in, -8.5in);
        
        %Lundi - Mardi
        \draw [gray!90] (3in, 1in) to (3in, -8.5in);

        %Mardi - Mercredi
        \draw [gray!90] (4.5in, 1in) to (4.5in, -8.5in);
        
        %Ligne séparatrice pour le nom des jours de la semaine
        \draw [gray!90] (1.5in, -\hauteurLigneJourSemaine in) to (6in, -\hauteurLigneJourSemaine in);
        

        

        
        \foreach \y [count=\counterOriginal, evaluate=\counterOriginal as \counter using ((\counterOriginal)-0.9)/1.79)] in \bulletCount {
            %Affiche la date du jour jusqu'au nombre de jour dans le mois.
            \ifdim\counterOriginal < \dayCountPlusOne
                %\fill (-0.18, -\counter) node (A) [day node] {\texttt{\y}};
            \fi

            %\ifdim\counterOriginal pt=\dayCountPlusOne pt
               % \fill (2.5, -\counter) node (A) [task node] {\text{\firstCat}}; 
              %  \fill (6.95, -\counter) node (A) [task node] {\text{\secondCat}}; 
             %   \fill (11.35, -\counter) node (A) [task node] {\text{\thirdCat}}; 
            %\else
               % \monthTaskDot{0.34, -\counter};
              %  \monthTaskDot{4.79, -\counter};
             %   \monthTaskDot{9.2, -\counter};
            %\fi
        }
    
        %\ifnum\daysCount=28
        %    \draw [loosely dashed] (0, -30.31) to (90,-30.31);  
       % \fi
       % \ifnum\daysCount=29
        %    \draw [loosely dashed] (0, -31.29) to (90,-31.29);  
       % \fi
       % \ifnum\daysCount=30
       %     \draw [loosely dashed] (-1, -16.53) to (15, -16.53);  
       % \fi
       % \ifnum\daysCount=31
       %     \draw [loosely dashed] (0, -33.25) to (90,-33.25);  
       % \fi


        
    \end{tikzpicture}
    
    
\end{document}